%% === WELCOME TO THE SOURCE CODE. HERE ARE SOME OF YANIV's EXTRA THOUGHTS/DISCUSSIONS/FAILED ATTEMPTS. ===
% \begin{enumerate}
%     \item Some clarification to Item-\ref{item_subclass} under \emph{LP Optimality Questions}: paths are of interest because they are the simplest tree topology without branching, and in their case the STTs are in fact BSTs. Edge-diameter of $3$ is interesting because it is plausible that the LP polytope in this case is still integer, as in the case of stars. From Example~\ref{example_non_int_vertex_dominated_in_D_space} we know that the path with $5$ nodes has a non-integer LP, and by Theorem~\ref{theorem_subgraph_with_frac_vertex_implies_frac_vertex_extension} the non-integrality is true for any other tree with edge-diameter of $\ge 4$. In particular, this non-integrality also holds in $D$-space for the long-star topology (whose edge-diameter is also $4$).
    
%     \item Root rounding gives a $2$-approximation because each rooting of a new node adds (at most) $\frac{1}{2}$ to the depth of other nodes, while each depth is already at least $\frac{1}{2}$ with respect to the current subtree, so overall by induction the depth of $v$ at most doubles in the process.
    
%     One hope to tighten the analysis is to show that ``enough'' nodes do not increase in depth in ``enough'' of the steps, for some notion of ``enough''. For example, if we could show that we can choose a root such that at most $O(1)$ nodes gain the extra depth, perhaps we could average these extras such that no node suffers too much. However, even checking on the small graphs, we see an example with $n=8$ where regardless of the root we choose (in the first rooting), $3$ out of $7$ of the nodes gain an extra $\frac{1}{2}$ of depth. This does not rule out $O(1)$, but it can pessimistically be interpreted to suggest that many of the nodes could gain depth in any such rounding.
    
%     We note that even if many nodes gain depth, but not too many, and if it is possible to average the gain in depth across the nodes, the approximation ratio may be improved. For example, if at most half of the nodes gain depth, and if we can average it nicely among them, then every node will gain roughly $\frac{1}{2}$ extra depth per $1$ original depth (compared to $\frac{1}{2}$). However it is not clear how to ensure an even distribution of the gain of depth.

%     \item Can we improve the approximation ratio by introducing an additive term? Property~\ref{property_tree_2_approx} is derived by showing how to take the LP optimum with some depth vector $D$, and derive an STT whose depths vector $D'$ satisfies $\forall i: D'_i \le 2D_i$, which yields a $2$-approximation. Can we improve the multiplicative constant at the expense of adding an additive term, for example $D''$ such that $\forall i: D''_i \le c D_i + k$ for a constant $c < 2$ and a ``useful'' expression $k$? Ideally, $k$ would be a constant, but even if not it can yield a class of cases where the approximation ratio is improved. As a concrete example, if we can generally prove $c=1$ and $k=O(\log \log n)$, then in the case where the average query (of OPT) costs $\Omega(\log n)$, we get an almost-$1$ approximation.

%     \item A note on the ``peeling'' in Item~\ref{item_peeling} of \emph{Improved LP Approximations}: If such peeling can work, there is likely a way to ensure that the decomposition does not have too many STTs. It could be done by peeling STTs up to some threshold, to get some point $Q'$ which is close enough to $Q$. As a concrete suggestion, we can work with integers by taking the matrix $X$ and define $X' \equiv \ceil{k \cdot X}$ for some large $k$, and then peel STTs from $X'$ (each STT will subtracts $1$ from some entries of $X'$) until the maximum column sum is at most some small value $k'$. By taking the best STT in such a decomposition, denote it $T$, we would get that $T$ is optimal up to an additive constant $k'$.

%     \item An issue with iterative rounding is that one must be able to define a smaller LP to solve when iterating. If we determine and fix the root each time, we indeed get a smaller problem (for each connected component as a sub-STT). However, usually the LP solution does not give us a clean root. Other ways to get a smaller problems are contractions of nodes or edges to get a smaller topology, however it is not clear how to apply such contractions correctly. Note that simply fixing the value of some $X_{ij}$ based on a solution that we found does not help us reduce the size of the LP, because we would still get the same LP, and same solution, in the next iteration.

%     \item The objective of the dual LP (Definition~\ref{definition_dual_LP_program}) is an unweighted sum of $R_{ij}$ variables. Based on the capping (in)equality we can substitute each of them by some pair $Q_{ikj}+Q_{jki}$, or even by an equal representation of each pair $\frac{\sum_{k \in (i \leftrightsquigarrow j)}{(Q_{ikj}+Q_{jki})}}{|(i \leftrightsquigarrow j)|}$. Can this alternative presentation help us?

%     \item Some more regarding Quadratic-programming (QP): QP is a generalization of linear programming, in which the constraints are still linear, but the objective function over a vector of variables $x$ is to minimize $x^T Q x + c^T x$ where $Q$ is a matrix, and $c^T$ is a row vector. Note that if $Q=0$ then this is a linear program. One can solve problems in polynomial time if the cost function is convex, otherwise the problem may be NP-hard. The problem is convex if $Q$ is positive-definite.
    
%     In our case, it is appealing to use QP to penalize for quadratic terms of the form $X_{ij} \cdot X_{ji}$. We can take the penalty to be so large that hopefully we can guarantee that the optimal solution would satisfy that $X_{ij} \cdot X_{ji} = 0$ for every pair $i,j$. If we could do that, it may get us ``closer'' to getting an STT since any feasible solution that is induced by an STT, satisfies $X_{ij} \cdot X_{ji} = 0$ for every pair $i,j$ (because at least one node is not an ancestor of the other).
        
%     If we number the $X$ variables such that $X_{ij}$ and $X_{ji}$ go in pairs and follow each other, the matrix $Q$ has a blocks-structure where each block is of size $2 \times 2$, a matrix that is (up to a constant) $0$ on the main diagonal and $1$ on the other diagonal. Unfortunately, this block is not positive-definite, and as a result $Q$ is not either.
        
%     An alternative approach to use QP would be to penalize the cost when $X_{ij} - X_{ij}^2$ is large, to ``push'' towards optimum where $X_{ij}$ is either $0$ or $1$. However, then $Q$ is simply the identity matrix, negated, which again is not positive-definite.
    
%     Note that QP may yet be useful to refine the problem, it just depends whether we can refine the cost function with a useful positive-definite expression.
% \end{enumerate}
